\documentclass[11pt,a4paper]{article}
\usepackage[margin=25mm]{geometry}
\usepackage{amsmath,amssymb,amsthm}
\newtheorem{theorem}{Theorem}
\newtheorem{corollary}[theorem]{Corollary}
\newcommand{\Aut}{\operatorname{Aut}}
\newcommand{\Sym}{\operatorname{Sym}}
\newcommand{\N}{\operatorname{N}}
\newcommand{\C}{\operatorname{C}}
\title{When a permutation action detects abstract isomorphism}
\author{Human Compression project: L071 recovery}
\date{Working manuscript; distinct review pending}
\begin{document}\maketitle
Let $G$ be a finite group and let $\pi:G\hookrightarrow\Sym(\Omega)$
be a faithful permutation representation. An ordered generating tuple
$t=(t_0,\ldots,t_{r-1})$ of involutions defines a graph on $\Omega$:
join $x$ to $\pi(t_i)x$ with an edge of colour $i$ whenever these vertices
are distinct. Colours are retained, including parallel edges with
different colours. A fixed point of $\pi(t_i)$ is exactly a vertex
with no edge of colour $i$ incident to it. Thus the graph determines
every permutation $\pi(t_i)$, including its fixed points.

This note supplies the precise action hypothesis in Leemans--Mixer,
\emph{Algorithms for classifying regular polytopes with a fixed
automorphism group}, Theorem~4.1. The proof uses the elementary
classification of finite $G$-sets by their point stabilizers. That
classification, and the normalizer action below, are standard ideas;
no new general group-theoretic result is claimed.

\begin{theorem}\label{main}
For two ordered generating involution tuples $t,t'$ of $G$, their
coloured graphs are isomorphic if and only if some element of
$\N_{\Sym(\Omega)}(\pi(G))$ simultaneously conjugates
$\pi(t_i)$ to $\pi(t'_i)$ for every $i$. The automorphisms of $G$
induced in this way form the image of the homomorphism
\[
 \N_{\Sym(\Omega)}(\pi(G))\longrightarrow\Aut(G),
 \qquad f\longmapsto\bigl[g\mapsto\pi^{-1}(f\pi(g)f^{-1})\bigr],
\]
whose kernel is $\C_{\Sym(\Omega)}(\pi(G))$.

Write the orbit decomposition of $\Omega$ as
\[
       \Omega\cong\coprod_{[H]}m_{[H]}\,G/H,
\]
where $[H]$ ranges over conjugacy classes of subgroups of $G$ and
$m_{[H]}$ is a nonnegative integer. An automorphism $\alpha$ belongs
to the displayed image if and only if
\[
                   m_{[H]}=m_{[\alpha(H)]}\quad\hbox{for every }H.
                                                        \tag{1}
\]
In particular, if these multiplicities are invariant under every
abstract automorphism, coloured graph isomorphism is equivalent to
abstract equivalence of the generating tuples under $\Aut(G)$.
\end{theorem}
\begin{proof}
A colour-preserving graph isomorphism $f$ sends the unique neighbour
of $x$ of colour $i$, when present, to the corresponding neighbour of
$f(x)$; it also sends vertices with no such neighbour to vertices with
no such neighbour. Consequently
$f\pi(t_i)f^{-1}=\pi(t'_i)$ on every vertex. Since both tuples
generate $G$, this identity gives $f\pi(G)f^{-1}=\pi(G)$.
Conversely, simultaneous conjugation of these permutations preserves
exactly the edges and colours of the graphs. The assertion about the
homomorphism and its kernel follows directly from its formula and
the faithfulness of $\pi$.

Suppose $f$ induces $\alpha$. The relation
$f(gx)=\alpha(g)f(x)$ shows that the stabilizer of $f(x)$ is
$\alpha(G_x)$. Thus $f$ permutes the transitive orbits according
to $[H]\mapsto[\alpha(H)]$, proving the necessity of (1).
For sufficiency, match the copies of $G/H$ bijectively with the
copies of $G/\alpha(H)$ using (1), and on each matched pair use
\[
                         gH\longmapsto\alpha(g)\alpha(H).
\]
This is well defined and bijective, and has the required equivariance.
If the chosen representative of the target stabilizer class differs
from $\alpha(H)$, compose with the usual equivariant coset
isomorphism for conjugate subgroups. The maps on the disjoint orbits
assemble into a permutation $f$ inducing $\alpha$. Finally abstract
equivalence of tuples is exactly simultaneous application of an
automorphism to their entries; surjectivity of the displayed
homomorphism proves the last assertion.
\end{proof}

\begin{corollary}\label{regular}
For every finite $G$, its regular action detects abstract isomorphism
of all ordered generating involution tuples by coloured graph
isomorphism. More generally, any faithful finite $G$-set with an
$\Aut(G)$-invariant multiset of stabilizer conjugacy classes does so.
For string C-group tuples, the same statement detects isomorphism
of the associated regular abstract polytopes. Reversing the colour
order detects duality.
\end{corollary}
\begin{proof}
The regular action has one orbit with trivial stabilizer. Every
automorphism preserves that stabilizer class, so Theorem~\ref{main}
applies. Explicitly, on the underlying set $G$, the permutation
$f_\alpha:g\mapsto\alpha(g)$ conjugates each left multiplication
$L_h$ to $L_{\alpha(h)}$. The more general assertion is the last
part of the theorem. The marked-group correspondence for regular
abstract polytopes identifies their isomorphisms with automorphisms
carrying the ordered distinguished generators to one another; that
standard correspondence is retained from the common theory.
Reversal of the ordered generators is the definition of the dual
marked polytope, giving the final assertion.
\end{proof}

\paragraph{Why two unrelated embeddings do not suffice.}
Take $G=\langle a,b\rangle\le\Sym(6)$ with
$a=(1\ 2)$ and $b=(3\ 4)(5\ 6)$. This is a Klein four group.
Both $(a,b)$ and $(b,a)$ are string C-group tuples: their distinct
order-two subgroups meet trivially, which verifies every intersection
of subgroups generated by subsets. They are abstractly equivalent
under the automorphism interchanging $a$ and $b$.
Nevertheless their coloured graphs are not isomorphic, since colour
zero has one edge in the first graph and two in the second.
The abstract group $\Aut(G)\cong\Sym(3)$ also embeds in $\Sym(6)$.
Thus merely embedding both abstract groups in the same symmetric
group does not supply the needed compatibility. In this example the
normalizer cannot induce the automorphism interchanging $a$ and $b$,
because conjugation preserves cycle type.

The source's proof uses the stronger assumption that automorphisms
act on the embedded $G$ by permutation conjugation. Under that
intended hypothesis, Theorem~\ref{main} proves its graph criterion.
Its literal statement requires this qualification. In the reverse
implication a normalizing permutation induces an automorphism; it
need not itself belong to a separately embedded copy of $\Aut(G)$.
The centralizer kernel in Theorem~\ref{main} makes this distinction
explicit.

\paragraph{Use in the finite classifications.}
The natural action of $A_n$ for $n\ge7$ satisfies the criterion,
using the retained classical theorem that all its automorphisms are
induced by $S_n$. In other models it is enough to exhibit a subgroup
$A\le\Sym(\Omega)$ normalizing $\pi(G)$ and verify that
$|A|/|\C_A(\pi(G))|=|\Aut(G)|$. The image is then the full
automorphism group. This is the exact guard used by the accompanying
enumerations. The regular action is always available, while the
stabilizer-multiplicity criterion explains when a smaller action
suffices.

This correction and elementary reproof do not settle an open
conjecture. The numerical classifications and their full finite
outputs remain separate premises of the L071 recovery. This note
alone supplies no whole-original, formalization or compression credit.
\end{document}
